\documentclass{article}

\usepackage[margin=1in]{geometry}
\usepackage{graphicx}
\usepackage{amsmath}
\usepackage{amssymb}
\usepackage{booktabs}
\usepackage{siunitx}
\usepackage{natbib}
\usepackage{subcaption}

\newcommand{\keywords}[1]{\par\noindent\textbf{Keywords:} #1\par}
\newcommand{\hstar}{H^{\star}}

\title{Right-Censoring of the Predictability Horizon: Why
Persistence-Relative Skill in Daily PM$_{10}$ Forecasting May Never Cross Zero}
\author{Francisco Federico Garcia Crespi}
\date{April 2026}

\begin{document}

\maketitle

% ============================================================
\begin{abstract}
Multi-horizon forecast evaluations routinely report a maximum evaluated
horizon, $H_{\max}$, and readers often treat it, implicitly or explicitly, as
an estimate of the predictability horizon $\hstar$: the lead time beyond which
a model no longer beats a reference forecast.
This inference is unsound whenever skill at $H_{\max}$ is still positive,
because a positive endpoint is compatible with $\hstar > H_{\max}$ as much as
with $\hstar = H_{\max}$---a right-censoring problem that is rarely named as
such in the forecast-verification literature.
We study this problem in daily PM$_{10}$ forecasting at three stations in
eastern Spain spanning an industrial/suburban, an urban/residential, and a
rural remote background site, evaluated under rolling-origin protocols with
persistence as the reference forecast for $h=1,\dots,7$ days.
At all three stations, and for the best of three candidate models (direct
gradient-boosted trees, direct ridge regression, sparse-origin SARIMA), skill
relative to persistence is positive at every evaluated horizon and shows no
sign of approaching zero by $h=7$.
We show that this pattern is not evidence of overfitting: skill tracks a
theoretical linear-predictability ceiling derived purely from the empirical
autocorrelation function (a generalized Yule--Walker bound using two weeks of
lags), and a block bootstrap over folds finds no statistically robust
exceedance of that ceiling in any of the 21 station-horizon cells examined.
We further show, algebraically, that this outcome is expected rather than
anomalous: under a persistence baseline, the skill of the best linear
predictor of a stationary, mean-reverting series is monotonically increasing
in $h$ and converges to one half as $h\to\infty$, never crossing zero for any
finite horizon.
Consequently, $\hstar$ defined as a zero-crossing against persistence is not,
in general, a finite well-defined quantity for series of this type, and
studies that report $\hstar = H_{\max}$ as a terminal estimate are reporting a
lower bound dressed as a point estimate.
We recommend reporting the autocorrelation-implied linear ceiling alongside
any horizon-truncated skill curve as a low-cost, purely algebraic
verifiability check before a predictability horizon is claimed.
\end{abstract}

\keywords{PM$_{10}$ forecasting; predictability horizon; right-censoring;
rolling-origin evaluation; persistence baseline; autocorrelation; Yule--Walker;
forecast skill}

% ============================================================
\section{Introduction}

A recurring exercise in multi-horizon environmental forecasting is to report
skill relative to a reference forecast---typically persistence or climatology
---as a function of lead time $h$, up to some maximum evaluated horizon
$H_{\max}$, and to summarize the resulting curve by a single number: the
horizon at which skill stops being positive.
This quantity, which we denote $\hstar$, is treated in practice as an estimate
of the predictability horizon of the underlying process: the point beyond
which the reference forecast is as good as, or better than, the model
\citep{murphy1992}.

The inferential step from an evaluated curve to $\hstar$ is sound only when
skill is observed to cross zero within the evaluated range.
When it does not---when skill at $H_{\max}$ remains positive---the honest
conclusion is that $\hstar \ge H_{\max}$, not that $\hstar = H_{\max}$.
Treating the latter as established is a right-censoring error familiar from
survival analysis and reliability engineering, where an event observed not to
have occurred by the end of a study window is recorded as ``at least this
long,'' not as ``exactly this long.''
The parallel is exact here: $H_{\max}$ is a design choice of the evaluation
protocol, not a property of the data-generating process, and a positive skill
value at $H_{\max}$ says as much about the protocol as about predictability.

This is not a hypothetical concern.
In an informal audit of eleven research repositories built around PM$_{10}$
and related environmental predictability studies, we found $H_{\max}$
hard-coded at a fixed value (typically 6, 7, 24, or 48) in nine of them, with
skill still positive at $H_{\max}$ in every case we could verify empirically,
and with no repository implementing a check that would flag $\hstar =
H_{\max}$ as a possible artifact of the evaluation window rather than a
substantive finding.
One internal analysis report came close to naming the issue---noting that an
$\hstar$ estimate ``saturates'' at the horizon cap without extending beyond
it---but stopped short of treating this as a statistical censoring problem
with algebraic consequences for what can and cannot be claimed.

This paper asks whether the pattern is real and, if so, whether it is
expected.
We use daily PM$_{10}$ series at three stations with markedly different
exposure regimes and evaluate persistence-relative skill for the best of three
forecasting approaches under a leakage-free rolling-origin protocol.
We then ask a sharper question than ``does skill stay positive'': given the
autocorrelation structure of each series, \emph{how much} skill is a linear
forecaster entitled to at each horizon, and does the empirical skill exceed
that entitlement.
Answering this requires only the empirical autocorrelation function of the raw
series---no additional model training---because the minimum mean-squared error
achievable by any linear combination of past values is determined by the
autocovariance structure via the Yule--Walker equations
\citep{boxjenkins2015}.

\paragraph{Contributions.}
\begin{itemize}
  \item We document, with fold-level evidence, that persistence-relative
    skill remains positive through the full evaluated horizon range at three
    PM$_{10}$ stations of different type (industrial/suburban, urban, rural
    remote background), and we show this is a case of right-censoring, not an
    established predictability limit.
  \item We derive and apply a rigorous linear-predictability ceiling from the
    empirical autocorrelation function via a generalized Yule--Walker
    construction, and we show, with a block bootstrap over rolling-origin
    folds, that empirical skill does not exceed this ceiling in any
    statistically defensible sense across three stations, up to seven
    horizons, and up to three candidate models.
  \item We give a closed-form argument for why $\hstar$ against a persistence
    baseline need not be a finite quantity for stationary, mean-reverting
    series, reframing the applied question from ``where does skill cross
    zero'' to ``does skill track the autocorrelation-implied ceiling.''
\end{itemize}

% ============================================================
\section{Related Work and Positioning}

This manuscript sits at the intersection of two established but rarely
combined literatures: baseline-relative forecast verification, and classical
linear time-series theory.

\subsection{Predictability horizons and reference forecasts}

Reporting skill relative to persistence or climatology across a set of lead
times is standard practice in atmospheric and environmental verification
\citep{murphy1992,donnelly2015}, and the resulting curve is frequently used,
informally, to bound the useful range of a forecasting system
\citep{gneiting2011}.
Rolling-origin (walk-forward) aggregation is the appropriate protocol for this
exercise because it preserves temporal causality and exposes the evaluation to
a changing sequence of conditions, rather than to a single, possibly
unrepresentative split \citep{tashman2000}.
What is less commonly made explicit is that any such curve is observed only up
to a finite $H_{\max}$, and that a positive endpoint is a lower bound on the
true predictability horizon, not an estimate of it.
This is the same logical structure as right-censored survival data, though the
connection is, to our knowledge, not routinely drawn in the forecast
verification literature.

\subsection{Linear predictability ceilings}

Classical time-series theory provides an exact answer to how much a linear
forecaster can achieve at a given horizon, given only the autocovariance
structure of a stationary process.
For an autoregressive process of order one, the minimum mean-squared error of
the optimal $h$-step predictor follows directly from the lag-1 autocorrelation
\citep{boxjenkins2015}.
More generally, the minimum mean-squared error achievable by any linear
combination of the $p$ most recent observations is obtained by solving the
Yule--Walker equations for the relevant autocovariance matrix
\citep{boxjenkins2015}.
This construction requires no model fitting beyond estimating the
autocorrelation function itself, and it provides a ceiling that applies to
\emph{any} linear predictor using at most $p$ lags of the series, not only to
the specific model class under evaluation.
We are not aware of this ceiling being used as a verifiability check on
reported skill curves in the applied air-quality forecasting literature,
where its role is closer to that of a diagnostic sanity check than to a
competing forecasting method.

\subsection{Positioning}

We do not propose a new model, a new skill score, or a new baseline.
The contribution is a verifiability argument: given only the raw target series
and fold-level predictions that most rolling-origin studies already produce
as an intermediate artifact, one can test, with pure algebra on the
autocorrelation function and no additional training, whether a reported skill
curve is internally consistent with the predictability implied by the series
itself, and whether an apparent horizon at which skill remains positive is
evidence of a real limit or an artifact of where evaluation stopped.

% ============================================================
\section{Data and Study Sites}
\label{sec:data}

We use daily-mean PM$_{10}$ series from three monitoring stations in eastern
Spain, chosen to span markedly different exposure regimes while all being
subject to the same evaluation protocol (Table~\ref{tab:sites}).
Elx-Agroalimentari (Alicante province) is a suburban/industrial site at low
elevation; Valencia-Vivers (Valencia city) is an urban/residential site;
Zarra-EMEP (Valencia province) is a rural remote background station at
\SI{855}{\meter} elevation, part of the European Monitoring and Evaluation
Programme network.
All three series span 2017-01-01 to 2024-12-30 at daily resolution, with
4.0--19.6\% of calendar days missing (Table~\ref{tab:sites}); missing days
were treated as missing, not imputed, throughout the analysis reported here,
including the autocorrelation estimation described in
Section~\ref{sec:acf_method}.

\begin{table}[t]
\centering
\small
\caption{Study sites. VR here denotes the fraction of calendar days present
  (not missing) in the raw series.}
\label{tab:sites}
\begin{tabular}{@{}lllrrr@{}}
\toprule
Station & Province & Type & Altitude (m) & Days present & Missing (\%) \\
\midrule
Elx-Agroalimentari & Alicante & Suburban/Industrial & 44  & 2{,}350 & 19.6 \\
Valencia-Vivers     & Valencia & Urban/Residential   & 11  & 2{,}679 & 8.3 \\
Zarra-EMEP          & Valencia & Rural Remote (EMEP) & 855 & 2{,}804 & 4.0 \\
\bottomrule
\end{tabular}
\end{table}

% ============================================================
\section{Methodology}

\subsection{Rolling-origin evaluation and candidate models}

At each station, we evaluate persistence-relative skill under a leakage-free,
expanding-window rolling-origin protocol for horizons $h=1,\dots,7$ days.
Three candidate forecasting approaches are compared against a persistence
baseline ($\hat{y}_{t+h\mid t}=y_t$): a direct multi-horizon gradient-boosted
trees model, a direct multi-horizon ridge regression model, and a sparse-origin
SARIMA model with fixed order $(p,d,q)(P,D,Q)_s=(1,0,1)(1,0,1)_7$, refit on an
expanding window approximately every two weeks rather than at every origin, as
is common practice for computationally expensive statistical models
\citep{boxjenkins2015}.
Two additional weak references---a seasonal (weekly) persistence model and an
STL-decomposition-plus-ridge model---are reported for context only and are
excluded from model selection, since neither is a serious forecasting
candidate at these sites.
For each station and horizon, we report the skill of whichever of the three
candidate models achieves the lowest mean squared error, matched exactly on
the rolling-origin fold identifier against the persistence baseline evaluated
on the same folds, so that models with different fold coverage (notably
sparse-origin SARIMA, with 99--172 folds against 1{,}215--1{,}479 for the
daily-evaluated models) are never compared against a persistence baseline
computed on a different sample.

\subsection{Persistence-relative skill}

For model $m$ and horizon $h$, skill relative to persistence is
\begin{equation}
  \mathrm{Skill}_m(h) = 1 - \frac{\mathrm{MSE}_m(h)}{\mathrm{MSE}_{p}(h)},
\end{equation}
where $\mathrm{MSE}_m(h)$ and $\mathrm{MSE}_p(h)$ are the model's and
persistence's mean squared errors over the matched fold set at horizon $h$.
Positive values indicate lower error than persistence.

\subsection{Autocorrelation estimation on incomplete calendars}
\label{sec:acf_method}

Because 4--20\% of calendar days are missing at these sites, autocorrelation
must be estimated with care.
Compacting the series by discarding missing days before computing lagged
correlations silently converts a nominal lag of $k$ days into a variable
calendar gap---sometimes substantially larger than $k$---which biases the
estimated autocorrelation function downward.
We instead reindex each series to the full daily calendar, leaving missing
days as \texttt{NaN}, and estimate the autocovariance at lag $k$ using
pairwise-complete observations around the series' global mean:
\begin{equation}
  \hat\gamma(k) = \frac{1}{|\mathcal{P}_k|}
  \sum_{(t,t+k)\in\mathcal{P}_k} (x_t-\bar{x})(x_{t+k}-\bar{x}),
\end{equation}
where $\mathcal{P}_k$ is the set of index pairs $k$ calendar days apart with
both observations present, and $\bar{x}$ is the sample mean over all present
observations.
The lag-$k$ autocorrelation is $\hat\rho(k)=\hat\gamma(k)/\hat\gamma(0)$, and
$\phi \equiv \hat\rho(1)$.

\subsection{Three predictability ceilings}
\label{sec:bounds}

We compare empirical skill against three ceilings of increasing rigor, all
computed algebraically from $\hat\rho(\cdot)$ with no model fitting.

\textbf{AR(1) ceiling.}
If the series were a pure first-order autoregressive process with lag-1
autocorrelation $\phi$, the minimum mean-squared error of the optimal $h$-step
predictor is $\gamma(0)(1-\phi^{2h})$, while the mean-squared error of
persistence under the same process is $2\gamma(0)(1-\phi^{h})$.
The implied skill ceiling is
\begin{equation}
  \mathrm{Skill}_{\mathrm{AR1}}(h) = 1 -
  \frac{1-\phi^{2h}}{2(1-\phi^{h})} = \frac{1-\phi^h}{2}.
  \label{eq:ar1}
\end{equation}
This is a restrictive special case: it is exact only if the series is
genuinely AR(1), and it will understate achievable skill whenever the series
has structure beyond a single lag, such as weekly periodicity.

\textbf{Hybrid ceiling.}
A tempting shortcut is to keep the AR(1) numerator in Equation~\ref{eq:ar1}
but replace the AR(1)-implied persistence denominator with the true empirical
one, $2(1-\hat\rho(h))$:
\begin{equation}
  \mathrm{Skill}_{\mathrm{hybrid}}(h) = 1 - \frac{1-\phi^{2h}}{2(1-\hat\rho(h))}.
  \label{eq:hybrid}
\end{equation}
We include this construction because it appears, in one form or another, in
practice as an intermediate step, and we show in Section~\ref{sec:results}
that it is internally inconsistent---the numerator still assumes AR(1)
dynamics that the denominator's own empirical $\hat\rho(h)$ contradicts---and
should not be used as a ceiling.

\textbf{Linear ceiling (generalized Yule--Walker).}
The rigorous construction replaces the AR(1) numerator with the true minimum
mean-squared error achievable by the best linear combination of the $p$ most
recent lags, obtained by solving the Yule--Walker equations for the $p\times
p$ Toeplitz correlation matrix $\mathrm{P}$ with entries $\mathrm{P}_{ij}=
\hat\rho(|i-j|)$:
\begin{equation}
  E_{\mathrm{opt}}(h) = 1 - \boldsymbol{\rho}_h^\top \mathrm{P}^{-1}
  \boldsymbol{\rho}_h, \qquad
  [\boldsymbol{\rho}_h]_i = \hat\rho(h+i),\ i=0,\dots,p-1,
  \label{eq:yw}
\end{equation}
\begin{equation}
  \mathrm{Skill}_{\mathrm{lin}}(h) = 1 -
  \frac{E_{\mathrm{opt}}(h)}{2(1-\hat\rho(h))}.
  \label{eq:linbound}
\end{equation}
We use $p=14$ (two weeks), sufficient to capture any weekly cycle in the daily
series without imposing it a priori.
Equation~\ref{eq:yw} is the normalized minimum mean-squared error of
predicting $x_{t+h}$ from $(x_t,\dots,x_{t-p+1})$ using the best linear
combination consistent with the estimated autocorrelation structure; it is
not specific to any parametric model class, and any predictor restricted to
linear combinations of these $p$ lags is bounded above by
Equation~\ref{eq:linbound}, whether or not it is called AR, ARIMA, or linear
regression.
This is the only one of the three ceilings that is a genuine upper bound for
linear forecasters; Equations~\ref{eq:ar1} and~\ref{eq:hybrid} are reported
for comparison, not as claims.

\subsection{Robustness of ceiling exceedances}
\label{sec:bootstrap_method}

Where empirical skill exceeds $\mathrm{Skill}_{\mathrm{lin}}(h)$ at a given
station-horizon-model cell, we assess whether the exceedance is statistically
distinguishable from sampling noise using a non-parametric bootstrap: 5{,}000
resamples with replacement of the matched fold set, recomputing skill on each
resample, and a 95\% percentile interval.
We flag an exceedance as robust only if the entire interval lies above the
ceiling.
Two caveats apply to this procedure and are stated explicitly wherever
relevant: origins evaluated daily (the gradient-boosted trees and ridge
models, with $h$ up to 7) have heavily overlapping evaluation windows, so an
i.i.d.\ bootstrap over folds likely \emph{understates} true sampling
variability, making non-rejection of the ceiling a conservative conclusion for
these models; SARIMA's sparse-origin design (approximately biweekly, chosen
independently of any predictability consideration) yields near-independent
folds but small samples ($n\approx 99$--$172$), so wide intervals for SARIMA
reflect limited data, not necessarily model behavior.

% ============================================================
\section{Results}
\label{sec:results}

\subsection{Skill remains positive through $h=7$ at all three stations}

At every station and for the best-performing candidate model, persistence-
relative skill is positive at every evaluated horizon and shows no sign of
approaching zero by $h=7$ (Table~\ref{tab:main_results}).
Best-model skill at $h=7$ ranges from 0.370 (Zarra-EMEP, ridge) to 0.470
(Elx-Agroalimentari, ridge), with no station showing a downward trend
suggestive of an imminent crossing.
Under the evaluation protocols used to generate these results, $H_{\max}=7$
is a hard limit enforced in code, not a value chosen after observing where
skill happened to cross zero; extending it requires modifying the
rolling-origin fold generator, not merely a configuration file.

\begin{table}[t]
\centering
\small
\setlength{\tabcolsep}{4pt}
\caption{Persistence-relative skill of the best candidate model at each
  horizon, against the three theoretical ceilings. $\hat\rho(1)=\phi$ is the
  lag-1 autocorrelation of the reindexed calendar series.}
\label{tab:main_results}
\begin{tabular}{@{}clrrrrrl@{}}
\toprule
Station & $h$ & $\hat\rho(h)$ & $\mathrm{Skill}_{\mathrm{AR1}}$ &
  $\mathrm{Skill}_{\mathrm{hybrid}}$ & $\mathrm{Skill}_{\mathrm{lin}}$ &
  Skill (best) & Model \\
\midrule
Elx-Agroalimentari & 1 & 0.525 & 0.238 & 0.238 & 0.252 & 0.202 & ridge \\
($\phi=0.511$)     & 2 & 0.208 & 0.363 & 0.417 & 0.403 & 0.402 & ridge \\
                    & 3 & 0.112 & 0.428 & 0.449 & 0.451 & 0.449 & ridge \\
                    & 4 & 0.073 & 0.462 & 0.464 & 0.470 & 0.469 & ridge \\
                    & 5 & 0.046 & 0.480 & 0.477 & 0.484 & 0.479 & ridge \\
                    & 6 & 0.051 & 0.490 & 0.474 & 0.481 & 0.482 & ridge \\
                    & 7 & 0.028 & 0.495 & 0.486 & 0.492 & 0.470 & ridge \\
\addlinespace
Valencia-Vivers     & 1 & 0.606 & 0.197 & 0.197 & 0.216 & 0.182 & ridge \\
($\phi=0.614$)      & 2 & 0.376 & 0.317 & 0.307 & 0.338 & 0.323 & hgb \\
                    & 3 & 0.294 & 0.389 & 0.327 & 0.377 & 0.366 & ridge \\
                    & 4 & 0.246 & 0.433 & 0.349 & 0.400 & 0.398 & ridge \\
                    & 5 & 0.199 & 0.459 & 0.380 & 0.424 & 0.425 & ridge \\
                    & 6 & 0.195 & 0.475 & 0.380 & 0.423 & 0.428 & ridge \\
                    & 7 & 0.188 & 0.485 & 0.385 & 0.425 & 0.442 & ridge \\
\addlinespace
Zarra-EMEP          & 1 & 0.628 & 0.186 & 0.186 & 0.227 & 0.172 & sarima \\
($\phi=0.643$)      & 2 & 0.317 & 0.303 & 0.382 & 0.373 & 0.351 & hgb \\
                    & 3 & 0.213 & 0.376 & 0.404 & 0.427 & 0.485 & sarima \\
                    & 4 & 0.184 & 0.422 & 0.402 & 0.441 & 0.461 & sarima \\
                    & 5 & 0.172 & 0.451 & 0.402 & 0.447 & 0.563 & sarima \\
                    & 6 & 0.160 & 0.469 & 0.407 & 0.453 & 0.511 & sarima \\
                    & 7 & 0.136 & 0.481 & 0.422 & 0.466 & 0.370 & ridge \\
\bottomrule
\end{tabular}
\end{table}

\subsection{Empirical skill tracks the linear ceiling, not the AR(1) ceiling}

The AR(1) ceiling is exceeded by the best model in 9 of 21 station-horizon
cells (all at $h\ge2$), a pattern that is expected rather than concerning: all
three sites show a partial autocorrelation bump around $h=2$--$4$ relative to
a pure geometric decay, consistent with sub-weekly structure that AR(1) cannot
represent by construction.
This is precisely the motivation for the linear ceiling of
Equation~\ref{eq:linbound}, which uses two weeks of lags rather than one:
across the same 21 cells, $\mathrm{Skill}_{\mathrm{lin}}(h)$ is at or above
$\mathrm{Skill}_{\mathrm{AR1}}(h)$ in all but two cells, and by construction
can only be tighter.

The hybrid ceiling of Equation~\ref{eq:hybrid} is not a reliable stand-in for
the linear ceiling: it differs from $\mathrm{Skill}_{\mathrm{lin}}$ by
0.028 on average and by up to 0.051 (Valencia-Vivers, $h=4$), with a
systematic downward bias at Valencia-Vivers and Zarra-EMEP of 0.02--0.05 at
most horizons.
Because it retains the AR(1) assumption in its numerator while using the true
$\hat\rho(h)$ in its denominator, it understates genuine linear predictability
whenever the series has structure beyond lag 1---exactly the sites where it
matters most.
We do not recommend its use as a predictability ceiling.

\subsection{No robust exceedance of the linear ceiling}

The best model's skill exceeds $\mathrm{Skill}_{\mathrm{lin}}(h)$ in 8 of the
21 cells (Table~\ref{tab:main_results}), by margins ranging from a negligible
$+0.001$ (Valencia-Vivers, $h=5$) to $+0.117$ (Zarra-EMEP, $h=5$, SARIMA).
The bootstrap procedure of Section~\ref{sec:bootstrap_method} finds none of
these eight exceedances to be statistically robust: in every case, the 95\%
percentile interval spans the ceiling (Table~\ref{tab:bootstrap}).
The four ridge exceedances (all $\le+0.017$, over 1{,}224--1{,}446 folds) sit
well inside intervals of width 0.15--0.31, and, as noted in
Section~\ref{sec:bootstrap_method}, the true interval for these
heavily-overlapping daily folds is likely wider still.
The four SARIMA exceedances at Zarra-EMEP, including the largest observed
margin, occur over only 99--104 near-independent folds and carry intervals of
width 0.56--0.63---wide enough that the point estimate provides essentially no
evidence against the ceiling.

\begin{table}[t]
\centering
\small
\caption{Bootstrap assessment (5{,}000 resamples, 95\% percentile interval) of
  the eight station-horizon-model cells where point-estimate skill exceeded
  $\mathrm{Skill}_{\mathrm{lin}}(h)$. None is a robust exceedance.}
\label{tab:bootstrap}
\begin{tabular}{@{}llrrrrr@{}}
\toprule
Station & Model & $h$ & $N$ folds & Skill & 95\% CI & $\mathrm{Skill}_{\mathrm{lin}}$ \\
\midrule
Elx-Agroalimentari & ridge  & 6 & 1{,}224 & 0.482 & [0.323, 0.638] & 0.481 \\
Valencia-Vivers     & ridge  & 5 & 1{,}446 & 0.425 & [0.349, 0.501] & 0.424 \\
Valencia-Vivers     & ridge  & 6 & 1{,}442 & 0.428 & [0.355, 0.498] & 0.423 \\
Valencia-Vivers     & ridge  & 7 & 1{,}434 & 0.442 & [0.359, 0.519] & 0.425 \\
Zarra-EMEP          & sarima & 3 &   104   & 0.485 & [0.153, 0.716] & 0.427 \\
Zarra-EMEP          & sarima & 4 &   101   & 0.461 & [0.131, 0.719] & 0.441 \\
Zarra-EMEP          & sarima & 5 &   101   & 0.563 & [0.140, 0.753] & 0.447 \\
Zarra-EMEP          & sarima & 6 &    99   & 0.511 & [0.132, 0.743] & 0.453 \\
\bottomrule
\end{tabular}
\end{table}

Taken together, empirical skill tracks $\mathrm{Skill}_{\mathrm{lin}}(h)$
closely and without statistically defensible exceedance across three sites of
different type, seven horizons, and up to three candidate models.
This is the internal-consistency check that a positive skill value at
$H_{\max}$ should always be able to pass before it is interpreted
substantively.

% ============================================================
\section{Discussion}

\subsection{Why $\hstar$ against persistence may not be finite}

Equation~\ref{eq:ar1} makes the underlying mechanism explicit.
For any stationary process with $0<\phi<1$,
$\mathrm{Skill}_{\mathrm{AR1}}(h)=(1-\phi^h)/2$ is strictly increasing in $h$
and converges to $\nicefrac{1}{2}$ as $h\to\infty$, without ever reaching
zero.
The mechanism is not specific to AR(1): it follows from persistence's error
growing without bound as $h$ increases (toward $2\gamma(0)$, since $\rho(h)\to
0$ for any mixing stationary process), while the best predictor's error is
bounded above by the unconditional variance $\gamma(0)$---the error of simply
forecasting the mean.
Under a persistence baseline and mean-reverting dynamics, skill is therefore
mathematically prevented from crossing zero at any finite horizon; it merely
approaches an asymptote from below.
This is not a peculiarity of PM$_{10}$: it holds for any stationary series
with positive autocorrelation evaluated against persistence, which describes
a large share of environmental time series reported in the literature.

The practical implication is a reframing of the applied question.
``Where does skill cross zero'' is not, in general, an answerable question
under a persistence baseline for series of this type; the informative
question is ``does skill track the autocorrelation-implied ceiling, or does
it exceed it.''
The former invites exactly the right-censoring error described in
Section~\ref{sec:data}: a study that stops evaluating at $H_{\max}=7$ because
skill is still positive, and reports $\hstar=7$, is reporting where the study
stopped, not where predictability ends.
The latter is answerable with the raw series and fold-level predictions
alone, requires no additional data collection or model retraining, and is
precisely what we report in Section~\ref{sec:results}.

\subsection{What a positive result at $H_{\max}$ should be reported as}

Where skill remains positive at $H_{\max}$ and tracks the linear ceiling, as
in all three stations studied here, the defensible statement is that the
evaluated models extract close to the full linear predictability available in
the series over the tested range, and that the study provides a lower bound,
$\hstar \ge H_{\max}$, rather than a point estimate.
Whether skill eventually declines beyond $H_{\max}$---for instance, once the
autocorrelation structure itself weakens enough, or once a model class exists
that legitimately exceeds the linear ceiling by exploiting non-linear or
exogenous information---is a separate empirical question that this study does
not resolve, and it cannot be resolved by re-examining the same evaluated
range more carefully.
It requires extending $H_{\max}$.

\subsection{Relation to variance-based diagnostics}

This paper's ceiling addresses a different failure mode from the variance-
retention diagnostics used elsewhere in this research program to detect
forecasts that achieve positive skill by collapsing toward a near-constant
value (a companion analysis in this same repository). The two diagnostics are
complementary and should be read together: variance retention asks whether a
skilled forecast is dynamically informative, while the linear ceiling
developed here asks whether a skilled forecast is even mathematically
surprising given the series' own autocorrelation. A forecast can pass one
check and fail the other, and a rolling-origin study aiming for a defensible
predictability claim should, in our view, report both.

% ============================================================
\section{Limitations}

The linear ceiling of Section~\ref{sec:bounds} bounds predictors restricted to
linear combinations of up to $p=14$ lags of the series itself; it is not a
formal bound on non-linear models (such as the gradient-boosted trees model
evaluated here) or on models using exogenous covariates, and its non-violation
by a non-linear model is an empirical observation in this dataset, not a
guarantee of the construction.
The bootstrap procedure treats each ceiling as fixed, since it is estimated
from approximately 2{,}900 days of raw series per station---substantially more
information than the 99--1{,}479 folds used to estimate model skill---but we
do not propagate the sampling uncertainty of $\hat\rho(\cdot)$ itself into the
ceiling, so reported intervals should be read as conservative with respect to
the ceiling.
The analysis is restricted to $h\le7$ days and to three stations in a single
region; whether the same tracking behavior holds at other lead times, other
pollutants, or sites with weaker or non-stationary autocorrelation structure
(e.g., strong trends or regime shifts) is not established here.
Finally, SARIMA's sparse-origin evaluation design, while methodologically
appropriate for a computationally expensive model refit on an expanding
window, yields small per-cell sample sizes at Zarra-EMEP (99--174 folds); the
resulting wide bootstrap intervals reflect this design choice and should not
be read as evidence either for or against SARIMA's skill claims at that site.

% ============================================================
\section{Conclusion}

Across three PM$_{10}$ stations of different type in eastern Spain,
persistence-relative skill remains positive through $h=7$ days for the best
of three candidate forecasting models, a pattern that, taken alone, is
routinely (mis)read as evidence of a predictability horizon at or beyond the
maximum evaluated lead time.
We show that this reading conflates a lower bound with a point estimate: skill
at $H_{\max}$ being positive establishes $\hstar\ge H_{\max}$, not
$\hstar=H_{\max}$, and, under a persistence baseline and mean-reverting
dynamics, $\hstar$ may not be a finite quantity at all.
Rather than leaving this as a caveat, we provide a concrete, low-cost check:
a linear-predictability ceiling computed algebraically from the empirical
autocorrelation function via a generalized Yule--Walker construction, requiring
no additional model training.
Empirical skill tracks this ceiling closely at all three stations, with no
statistically robust exceedance across 21 station-horizon cells, indicating
that the observed positive skill is consistent with---rather than anomalous
relative to---the predictability implied by each series' own autocorrelation
structure.
We recommend this ceiling, or an equivalent, be reported alongside any
horizon-truncated skill curve before a predictability horizon is claimed from
it.

% ============================================================
\section*{Data and Code Availability}

Raw daily PM$_{10}$ series, rolling-origin fold-level predictions, and the
scripts used to compute the autocorrelation function, the three theoretical
ceilings, and the bootstrap robustness check are available in the companion
repository at
\url{https://github.com/fedeg-umh-es/varret-pm10-paper}.
Model predictions for the direct gradient-boosted trees and ridge regression
candidates originate from the companion repository at
\url{https://github.com/fedeg-umh-es/p34-variance-retention-api}.
Raw PM$_{10}$ data were sourced from Spanish national and regional air-quality
monitoring networks; station coordinates and classifications are reported in
Table~\ref{tab:sites}.

\bibliographystyle{plainnat}
\bibliography{references}

\end{document}
